\documentclass[10pt]{article}
\usepackage[margin=0.75in]{geometry}
\usepackage{booktabs,tabularx,array,multirow}
\usepackage[T1]{fontenc}
\usepackage{lmodern}
\usepackage{enumitem}
\usepackage{xcolor,float}
\usepackage[font=small,skip=3pt]{caption}
\usepackage[hidelinks]{hyperref}
\setlength{\textfloatsep}{6pt}\setlength{\intextsep}{6pt}
\setlist{nosep,leftmargin=1.2em}
\setlength{\parskip}{3pt}
\setlength{\parindent}{0pt}
\newcolumntype{Y}{>{\raggedright\arraybackslash}X}
\newcommand{\good}[1]{\textbf{#1}}
\newcommand{\bad}[1]{\textcolor{red}{#1}}

\title{\vspace{-2.2em}CS839 HW2}
\author{YOUR NAME (aaneja4@wisc.edu)}
\date{\vspace{-1.5em}}

\begin{document}
\maketitle

%=====================================================================
\section{Setup}

\textbf{Domain.} Airline flight operations, extended from the domain we set up in Assignment 1. Since we want
to compare multi-agent systems with a single agent, we split the data into one CSV per airline and added a
separate kind of data, hourly weather at every airport. This gives natural ways to divide the work: by airline,
or by role (a flight analyst and a weather analyst). We generated the data synthetically, so we know the true
answer to every question. Every task must be answered by running Python on the files, so each answer is a
number or a list we can grade exactly.

\textbf{Data.} 6 US airlines, 7 airports (ATL, DEN, DFW, JFK, LAX, ORD, SEA), January--February 2025.

\begin{table}[H]
\centering\small
\begin{tabularx}{\linewidth}{l l Y}
\toprule
File & Size & What it holds \\
\midrule
\texttt{airlines/\{WN,DL,AA,UA,B6,AS\}.csv} & 1,944 flights & One file per airline: flight id, date, origin, dest, scheduled and actual departure, cancelled \\
\texttt{weather.csv} & 9,912 rows & Hourly weather code and visibility at each airport \\
\texttt{airports.csv} & 7 rows & Airport, city, time zone, UTC offset \\
\texttt{README.md}, \texttt{weather\_codes.md} & docs & Column meanings, ``delayed = more than 15 min'', the exact bad-weather rule \\
\bottomrule
\end{tabularx}
\end{table}

\textbf{Definitions (from the README).} \emph{Departure delay} = actual departure $-$ scheduled departure, in
minutes; early departures are negative. A flight is \emph{delayed} if its departure delay is more than 15 minutes.
\emph{Delay rate} = delayed flights / completed flights, as a percentage. Cancelled flights are left out of both
the top and the bottom.

\textbf{Traps.} Simple counts over these files would be easy for any agent, so we added four traps. Each
one is a piece of information that has to be carried from the docs to whichever agent does the work, and each
is a mistake that several agents can easily make the same way. All four are written in the README, so a careful
agent can avoid them:
\begin{itemize}
\item \textbf{UTC:} Alaska (AS) stores times in UTC; all other airlines use local time.
\item \textbf{Duplicates:} JetBlue (B6) has 12 repeated rows.
\item \textbf{Midnight:} only clock times are stored. A flight scheduled 23:40 that leaves 01:31 gives
  $-1329$ min by plain subtraction; the real delay is $+111$ (add one day). But a flight that leaves 6 min early
  really is $-6$; adding a day there (``over-wrap'') makes it look 1434 min late.
\item \textbf{Weather rule:} ``bad weather'' is an exact rule (TS, SN, +RA, or fog with visibility $<$ 1 mile),
  not ``any rain or fog''.
\end{itemize}

%=====================================================================
\section{Agent Systems}

All three systems use the same model (\texttt{gpt-5.6-luna}, OpenAI Responses API) and the same budget of
\textbf{20 model calls per task}. Every helper agent is one-shot: it gets one message, does its job, returns
a result, and is gone. Agents talk only through JSON that our Python code passes between them.

\textbf{Tools.} \texttt{run\_python} runs code in a sandbox folder that holds only the files that agent is
allowed to see. \texttt{submit\_answer} returns the final JSON answer. Every tool result also says how many
calls are left, and the last call is forced to be a submit.

\begin{table}[H]
\centering\small
\begin{tabularx}{\linewidth}{l l l Y l}
\toprule
System & Agent & Tools & What it sees & Calls \\
\midrule
Single & Solver & python, submit & All files, README, the question & 20 \\
\midrule
\multirow{2}{*}{A: Hierarchy} & Manager & none (writes JSON) & Question, README, file list with columns; later the worker reports & 2 \\
 & Workers (1--3) & python, submit & Only their files and the manager's instructions (not the question) & 18 shared \\
\midrule
\multirow{3}{*}{B: Debate} & Solvers (3) & python, submit & Same as single, each works alone & 18 shared \\
 & Critic & none (writes JSON) & The 3 solvers' code, results and notes, plus the docs & 1 \\
 & Aggregator & none (writes JSON) & The 3 answers and the critique & 1 \\
\bottomrule
\end{tabularx}
\end{table}

\begin{itemize}
\item \textbf{Single agent.} One agent with everything. It decides itself when to check its work and when to stop.
\item \textbf{A: Hierarchy.} The manager splits the question into up to 3 subtasks, picks the files for each
  worker and writes their instructions. Workers run one after another. The manager then combines their reports
  in the same conversation, so it still remembers its plan.
\item \textbf{B: Debate.} Three solvers answer the full question on their own. The critic looks for bugs without
  running code. The aggregator picks the final answer.
\end{itemize}

\textbf{Budget split.} Shared pools are split evenly, and calls an agent does not use roll over to the next one.
We added this after a pilot with 11 calls per task: team agents got only 3 calls each, ran out, and submitted
``unable to compute''.

%=====================================================================
\section{Task Suite}

Eight tasks, two per category. Each category is picked to favour a different kind of system, so we can see
when a team helps and when it does not.

\begin{table}[H]
\centering\small
\begin{tabularx}{\linewidth}{l l Y Y Y l}
\toprule
ID & Category & What it asks & Why we picked it & Where it can fail & We expect \\
\midrule
T1 & Separable & Feb delay rate per airline & Each airline is its own piece, easy to split & Midnight, UTC, duplicates & Single (cheaper) \\
T2 & Separable & Per airline: busiest route and worst-delay route & More work per airline, still splits cleanly & Midnight, duplicates & Single (cheaper) \\
T3 & Distinct roles & DL+AA+UA: \% of delayed flights that left in bad weather & Needs a flight expert and a weather expert & Loose weather rule, midnight & Hierarchy \\
T4 & Distinct roles & AS only: delay rate in bad vs.\ other weather & Hardest join: UTC + weather + delay & UTC, weather rule, midnight & Hierarchy \\
T5 & Integration & One delay rate for the whole network & Parts must be merged the right way & Averaging airline rates instead of pooling counts & Single / debate \\
T6 & Integration & Top 5 routes over all airlines & Parts must be merged the right way & Merging per-airline top-3 lists & Single / debate \\
T7 & Wasteful & B6 cancellations on Feb 14 & One filter; a team can only add cost & Nothing & Single \\
T8 & Shared mistake & Average departure delay, all flights & One easy-to-make bug everyone can share & Midnight (forget it or overdo it) & Debate (hope) \\
\bottomrule
\end{tabularx}
\caption{The task suite, why each task is there, and which system we expect to do best.}
\end{table}

From this we wrote three hypotheses:
\textbf{H1:} hierarchy beats single on distinct-role tasks (T3, T4), because each worker gets a small, focused job.
\textbf{H2:} debate catches the shared mistake (T8), because the critic compares three independent solutions.
\textbf{H3:} on easy or separable tasks (T1, T2, T7) teams get the same score for 2--4$\times$ the cost.

\subsection{How we evaluate}
\begin{itemize}
\item \textbf{Score.} Graded by code against the known answer. Numbers must be within $\pm0.1$. Answers with
  several parts (6 airlines, 5 routes) get partial credit per part. 1.0 means fully right.
\item \textbf{Trap match.} We also compute the answer each trap would give. If a wrong answer equals one of them,
  we know which mistake the agent made.
\item \textbf{Cost.} Model calls, input/output tokens, and dollars (\$0.10 / \$0.50 per 1M tokens).
\item \textbf{Coordination.} Handoffs and characters passed between agents. For hierarchy: how many workers the
  manager used and which README caveats it passed down. For debate: how many different answers the solvers gave
  and whether the final answer matched the majority.
\item \textbf{Repeats.} T4 and T8 were the closest comparisons, so we ran them 3 times for every system.
\end{itemize}

%=====================================================================
\section{Results}

We first ran everything with the model's default reasoning. All three systems scored about 100\% (single 1.00,
hierarchy 1.00, debate 0.99), so there was nothing to compare. We then reran every agent with \textbf{low
reasoning effort}. The table below is from those runs.

\begin{table}[H]
\centering\small
\begin{tabularx}{\linewidth}{l ccc Y}
\toprule
Task & Single & Hierarchy & Debate & What happened and why \\
\midrule
T1 & 1.00 & 1.00 & 1.00 & One Python loop covers all 6 files, so splitting adds nothing. \\
T2 & 1.00 & 1.00 & 1.00 & Same. The only task where the manager used more than one worker. \\
T3 & 1.00 & 1.00 & 1.00 & The weather rule is written out exactly, so nobody got it wrong. \\
T4 (3 runs) & 0.67 & \good{1.00} & 0.67 & The manager turned every README caveat into a checklist for its worker. Single once over-wrapped. Debate once had all 3 solvers wrong. \\
T5 & 1.00 & 1.00 & 1.00 & The merge trap never fired: the manager used 1 worker, so nothing was merged. \\
T6 & 1.00 & 1.00 & 1.00 & Same. \\
T7 & 1.00 & 1.00 & 1.00 & Everyone right, but debate cost 5$\times$ more than single. \\
T8 (3 runs) & 0.67 & 0.67 & 0.67 & The same wrong numbers (1.0 min, 13.9 min) showed up in all systems. In debate the critic picked a wrong solver over a correct majority. \\
\midrule
Mean score & 0.83 & \good{0.92} & 0.83 & \\
Calls per task & 4.8 & 7.9 & 16.7 & \\
Cost per task (\$) & \good{.0020} & .0034 & .0075 & \\
Handoffs / chars passed & 0 / 0 & 2.2 / 4.4k & 3 / 12.2k & \\
\bottomrule
\end{tabularx}
\caption{Per-task score at low reasoning (1 run per cell unless marked).}
\end{table}

\textbf{Coordination.} The hierarchy manager almost never split the work: it used 1 worker on 7 of 8 tasks.
Its real job became rewriting the README as clear instructions. On every T4 run it passed all six caveats we
track (threshold, cancellations, UTC, duplicates, midnight, weather rule).

%=====================================================================
\section{Topology Surgery}

We made two structural changes, each aimed at a problem we saw in Section 4, and reran only the tasks it affects,
3 times each at low reasoning.
\begin{itemize}
\item \textbf{Hierarchy $\rightarrow$ forced split.} The manager kept using 1 worker, so the integration traps never
  got tested. We now require exactly 3 workers with 2 airlines each. Nothing else changes.
\item \textbf{Debate $\rightarrow$ majority vote.} The critic was making wrong calls, so we removed the critic and
  aggregator. The final answer is a plain vote over the 3 solvers (if all differ, take solver 1).
\end{itemize}

\begin{table}[H]
\centering\small
\begin{tabular}{l cc | l cc}
\toprule
Task & Hierarchy (Sec.\ 4) & Forced split & Task & Debate (Sec.\ 4) & Majority vote \\
\midrule
T2 & 1.00 & 0.86 & T1 & 1.00 & \bad{0.44} \\
T5 & 1.00 & \bad{0.33} & T2 & 1.00 & 1.00 \\
T6 & 1.00 & 0.93 & T4 & 0.67 & \good{1.00} \\
   &      &      & T8 & 0.67 & \good{1.00} \\
\midrule
Mean score & \good{1.00} & 0.71 & Mean score & 0.75 & \good{0.86} \\
Cost per task (\$) & \good{.0034} & .0059 & Cost per task (\$) & .0086 & \good{.0071} \\
\bottomrule
\end{tabular}
\caption{Surgery vs.\ the Section 4 result on the same tasks. Surgery columns are 3 runs per task.}
\end{table}

\textbf{Forced split: worse and costlier.} The manager obeyed every time, and it always merged raw counts
correctly, so the averaging trap still never fired. It failed in other ways:
\begin{itemize}
\item \textbf{One bad worker spoils the total} (T5, 2 of 3 runs). One worker over-wrapped and reported 55\% delays
  for its airlines (true is about 24\%). The manager added it in without noticing.
\item \textbf{Adding by hand} (T6). The manager cannot run code, added about 30 route counts in its head, and got one wrong.
\item \textbf{A missing piece} (T2). One worker returned nothing; the manager left 2 airlines blank rather than guess.
\end{itemize}

\textbf{Majority vote: better on T4 and T8, worse on T1.} Without the critic there is nobody to overrule a correct
majority, so T4 and T8 went to 3/3. But on T1 two solvers made the exact same over-wrap mistake, and the vote locked
it in. In the original debate the critic had rescued T1. A vote only works when mistakes are independent.

%=====================================================================
\section{Trajectories}

Three runs that show, step by step, how a team can help, hurt, or repeat a mistake. Full traces are in \texttt{traces/}.

\textbf{1. A team beat the single agent (T4, low reasoning).}
\begin{itemize}
\item \textbf{Single, run 0:} converted AS times from UTC and applied the weather rule correctly. But it added a day
  to \emph{every} negative delay, so flights that left a few minutes early counted as almost a day late. Score 0.
\item \textbf{Hierarchy:} the manager used one worker and wrote it a checklist: convert UTC, wrap midnight only for
  real overnight departures, delayed means more than 15 min, the exact bad-weather rule. The worker followed it and
  got 41.7\% (bad weather) vs.\ 6.8\% (other) on all 3 runs.
\item \textbf{Why:} the manager added no new data work. It helped because it turned the README into clear steps.
\end{itemize}

\textbf{2. Coordination hurt (debate T8, low reasoning, run 0).}
\begin{itemize}
\item Solvers 1 and 3 both got the right answer, 15.4 min. Solver 2 got 13.9 min (only half fixed the midnight wrap).
\item The critic cannot run code. Reading the code, it judged solver 2's method ``best supported''.
\item The aggregator followed the critic instead of the 2-of-3 majority. Final answer 13.9, score 0.
\item \textbf{Why:} the review step had more power than the solvers but less evidence. It overruled correct work.
\end{itemize}

\textbf{3. Agents repeated the same mistake (majority vote T1, low reasoning, run 0).}
\begin{itemize}
\item Solvers 2 and 3, working separately, produced \emph{identical} wrong answers (WN 53.6, DL 48.4, AA 48.8, \ldots),
  about 50\% delays for every airline, by over-wrapping early flights.
\item Solver 1 was right (WN 33.2, DL 20.2, \ldots). The vote went 2 to 1 for the wrong answer. Score 0.
\item \textbf{Why:} the solvers are the same model with the same prompt, so they make the same mistake. The same
  over-wrap also showed up in single T4 and in hierarchy workers. A vote cannot fix a mistake the majority shares.
\end{itemize}

%=====================================================================
\section{Takeaways}

\begin{enumerate}
\item \textbf{H1, hierarchy on distinct roles (T3, T4): partly true.} We expected it to win because each worker
  gets one job. It did win T4 (3/3 vs.\ 2/3), but not for that reason: it used one worker, not two experts. It won
  because the manager rewrote the README as a short checklist. When we forced it to really split, it got worse,
  because each extra worker is one more chance to make a mistake that nobody checks.
\item \textbf{H2, debate catches the shared mistake (T8): false.} We expected the critic to spot the midnight bug.
  Instead the critic, who cannot run code, twice picked a wrong answer over right ones. Removing it fixed T8 but
  let a shared wrong answer win on T1. Three solvers from the same model tend to make the \emph{same} mistake, so
  neither a critic nor a vote is a reliable check.
\item \textbf{H3, teams waste budget on easy tasks (T1, T2, T7): true.} Every system scored 1.0 on these. Hierarchy
  cost about 1.7$\times$ single on average and debate about 3.7$\times$, up to 5$\times$ on T7. A team only pays off
  when the extra agent adds real information (like the manager's checklist) or a check that can actually run code.
\end{enumerate}

\textbf{Limits.} Most cells are 1 run, so small gaps may be noise; only T4 and T8 have repeats. Total API cost for
all runs, including the pilot: about \$0.50.

\end{document}
